\documentclass[12pt]{article}
\usepackage{natbib}
\usepackage{pgf}
  \bibliographystyle{plainnat}
  \setcitestyle{citesep={,},aysep={}}
\usepackage[colorlinks,citecolor=blue,linkcolor=black,bookmarks=false,hypertexnames=true]{hyperref} 
\usepackage{url}
\usepackage[inline]{enumitem}
\usepackage{float}
\usepackage[bottom,perpage,multiple]{footmisc}
\usepackage{booktabs}
\usepackage{graphicx}
\usepackage{xcolor}
\usepackage{setspace}

\usepackage{enumitem}
\setlist{nosep}

  \setstretch{1.1}
\setlength{\bibsep}{0.0pt}
\tolerance=800

\usepackage{pgfplots}
%\DeclareUnicodeCharacter{2212}{-}
\usepgfplotslibrary{groupplots,dateplot}
\usetikzlibrary{patterns,shapes.arrows}
\pgfplotsset{compat=newest}

\usepackage{url}
\renewcommand*{\bibfont}{\footnotesize}

\title{
  \includegraphics[height=48pt]{logo}\\
  \textsc{Aibolit}:\\
  Style Checking\\
  Using Machine Learning}


\begin{document}

\maketitle

\begin{abstract}
Quality of code is crucial for the stability and robustness
of software systems. Defect detection at early stages of
software development lifecycle is the cheapest and the easiest
way to increase the quality of code. There are well-known
instruments such as static analyzers and linters, which are
actively used by programmers. However, none of them use Machine
Learning (ML) to detect defects more effectively. 
We managed to create such an instrument and empirically demonstrate 
its effectiveness\footnote{This is statement is yet to be confirmed}. 
\end{abstract}

\pagebreak

\section{Introduction}
% SPDX-FileCopyrightText: Copyright (c) 2019-2025 CQFN.org
% SPDX-License-Identifier: MIT

Insufficient software quality may result in increased development costs and
negatively affect customer satisfaction ~\citep{The_Economics_of_Software_Quality}.
\textit{Static code analysis} develops techniques to help detect software quality
issues prior to program execution. It has its practical application in various developer's tools. There are both open source\footnote{PMD: \url{http://pmd.sourceforge.net/}, Rubocop: \url{https://github.com/rubocop-hq/rubocop},
PHPCS: \url{https://github.com/squizlabs/PHP_CodeSniffer}
FindSecBugs: \url{https://find-sec-bugs.github.io/}, ESLint: \url{https://eslint.org/}, Checkstyle: \url{https://checkstyle.sourceforge.io/}.}
(PMD, Rubocop, PHPCS, FindSecBugs, ESLint, Checkstyle, to name a few) and commercial\footnote{IBM Security AppScan: \url{https://www.hcltechsw.com/wps/portal/products/appscan},
PVS-Studio: \url{https://www.viva64.com/en/pvs-studio/},
SonarQube: \url{https://www.sonarqube.org/},
Parasoft: \url{https://www.parasoft.com/}}
(IBM Security AppScan, PVS-Studio, SonarQube, Parasoft) static analyzers
on the market.

Static code analysis can be applied to improve an \textit{internal} and an
\textit{external} quality of software \citep{Ilyas2016StaticCA}. External
quality is related to defects encountered by the end user of the software
product. Within internal quality, two important subcategories are
\textit{functional quality} and \textit{maintainability}. Functional quality is
about code correctness and compliance with the functional software
specifications \citep{Farhan}. Code maintainability is about how easy it is to
analyze, modify and adapt given software \citep{Mohammadi2013AnAO}.

Functional quality aspects are typically quite susceptible to formal definition
and quantification.
% (\todo: examples!!).
Functional quality is also an essential
requirement in any domain of software development. On the other hand,
maintainability is a lot less straightforward to formally specify or quantify.
%\todo: refs.
Also, in certain applications it appears less important than
functional correctness, although in business domain it is recognized as an
essential property.
% (\todo: ref).
As a result, there are currently a lot more
research and practical tools addressing functional quality aspects of code than
maintainability \citep{Overview_Static_Code_Analysis_in_Software_Development}.
Another aspect of static analysis tools that may have hindered their application
to maintainability, is that they are predominantly rule-based. Since there has
not yet been a consensus on how to formalize maintainability, it is challenging
to devise a set of formal rules to detect it.

We designed our new tool Aibolit to help developers identify patterns in their
code  that may cause maintainability issues. It is a next generation static
analysis tool that uses a machine learning (ML) model as an underlying quality
prediction mechanism. From the perspective of ML, our product is a recommender
system. For a given class file, it gives suggestions to the developer to alter
their code. The recommendations come in the form of \textit{code patterns} that
are detected in the code and advised to be removed.

Our choice to design Aibolit as a ML-based system alleviates some important
shortcomings of rule-based static analyzers. By design, ML algorithms capture
statistical relations in the external world (data). Therefore, they can be a
good way to model imprecisely and subjectively defined properties of code, such
as its maintainability. Moreover, rule-based system are known to not scale well
to the diversity of empirically observed cases, and they tend to get very hard
to extend and maintain \citep{LenatFeigenbaum1987}. The ML
approach does not require
manual system adaptation as new observations or new features (patterns) come
along. In fact, Aibolit provides an easy way for developers to integrate a code
pattern of their liking into the recommender system and to analyze the pattern`s
impact on code quality.






% This document explains how Aibolit works  and what makes it novel and
% different from other static analyzers (Section~\ref{sec:method}).
% We further show how Aibolit can be extended
% with custom patterns and metrics. We discuss the current shortcomings and the way
% they can be addressed in future work (Section~\ref{sec:conclusion}).


\section{Related Work}
\label{sec:related}
There are many definitions of a defect.  
\citet{5989519} says that defect is ``a fault, bug, inaccuracy or lack of expected 
functionality in a project artifact.'' 
\citet{Assurance} says that it is ``a problem 
(synonym of fault) which, if not corrected, 
could cause an application to either fail or to produce incorrect results.''

Defects and software quality are directly related. There are multiple studies 
of defect types, their impact, complexity, 
root causse and other characteristics,~\citet[e.g.][]{10.1145/69605.2085, 
10.5555/256664.256773, 10.1145/390016.808455, Glass1981PersistentSE, 
10.1145/1353535.1346323, 10.1007/s10664-013-9258-8, catolino2019bugs}.
There are a few common preventive ways to deal with defects, like 
static code analyzers, testing software, or peer code review. 

Static code analyzer 
is a tool that helps find defects before a program is executed. 
Such an analyzer inspects various program representations, for example 
Abstract Syntax Tree (AST), Control Flow Graphs (CFG), or
Program Dependency Graph (PDG), 
and search for handcrafted defect patterns. 
These tools are popular among developers and are often embedded into
Integrated Development Environments (IDE) such as IntelliJ IDEA or NetBeans. 
There are more than 40 static analyzers currently on the market, including
very famouns open source projects, such as
PMD, Rubocop, PHPCS, Sparse, CLion, cpplint, FindSecBugs, ESLint, and Checkstyle.
There are also many commercial tools, like 
IBM Security AppScan, PVS-Studio, SonarQube, and Parasoft.
Software companies like Google and Facebook have their own open source
static analyzers: Error Prone~\citep{10.1109/SCAM.2012.28} and 
Infer~\citep{10.1007/978-3-319-17524-9_1} respectively. 
Static code analyzer may vary in supported languages, 
supported defect types, and their integration workflow.

Usually static code analyzers are rule-based. 
In order to detect a new kind of defect their developers have to
design a new pattern. To address this extensibility problem 
there were attempts to learn pattern
from data, as explained by~\citet{bielik2016learning, wang2019learning}.

Recent empirical study by~\citet{10.1145/3238147.3238213} demonstrates that 
state-of-the-art code analyzers miss more than 90\% of defects. Most of those missed 
defects are inconsistencies with the specification or programmer's intent. 
In order to catch such defects it is necessary to reason about 
possible behavior depending on the input data, 
which is difficult or impossible for the rule-based approach. 
Such defects are known as \emph{semantic} defects. 
It has been demonstrated by~\citet{10.1007/s10664-013-9258-8} 
that semantic defects are the dominant root cause for the majoirity of security issues, 
when attacker may get unauthorized access to some resources.

Another issue of most static analyzers is their high 
percent of ``false-positives,'' when they find a defect, which
in reality is not a defect. The amount of these false signals grows 
when the size and complexity of the project increases. 
This leads to developers loosing trust to the tool and stop
using it. 

The efficiency of static analyzers is yet another problem. It was
empirically shown by~\citet{10.1145/3188720} that 
a better performance metric for a static analyzer 
is the amount of fixed, rather than found defects. 
Static analyzers must provide the right information 
at the right time doing everything possible to not annoy 
software developers. 

Usually, static code analyzers use fixed in time code versions as their
input. However, one can also use the history of code changes and 
the information from the Issue Tracking System (ITS) 
to enhance the quality of prediction, as done by~\citet{Gupta2018IntelligentCR, kapur2018estimating}.

There are studies trying to use ML in order to detect defects. 
For example, \citet{Dam2018ADT} first trained vector representations of 
an AST in an unsupervised manner and then used it as a feature vector 
to train the binary classifier. 
\citet{kapur2018estimating} 
combined the information extracted from the code programming style and ITS 
and built a predictor to estimate the defectiveness of an input source code 
file. Using an idea that the names of identifiers (variables, classes, functions) 
convey useful information, which might be used to understand programmer's 
intent,
\citet{Pradel2018DeepBugsAL} proposed a method that first learned 
vector representation of identifiers 
and obtained a fixed length vector for a code snippet to train a binary classifier 
with feed-forward neural network. 
\citet{vasic2019neural} used pointer-network to do joint prediction of 
both the location and the possible fix for variable misuse bugs. 
\citet{briem2019using} used attention-based neural network  to model binary 
classifier to detect off-by-one defects.


\section{Categories of Methods}
\label{sec:categories}
Even though detecting defects is an important task in any software project,
existing studies differ in their definition of the parameters and objectives
of this task,
partially due to the ambiguity of the definition of the software defect.
Among others, there are many methods of defect detection, which utilize ML. 
Below we classify them from a few different perspectives, 
in order to justify the selection made later in this document.

\subsection{By the Origin of Data}

As suggested by~\citet{10.1007/s10462-017-9563-5}
the model can be built using the data 
\begin{enumerate*}[label=\arabic*)]
\item from the same release of a software module (\emph{intra-release}), 
\item from a few different releases of the same module (\emph{inter-release}), or 
\item from different software modules (\emph{cross-project prediction}). 
\end{enumerate*}
It seems that this classification was not possible earlier,
when analytical data was kept proprietary and was not available
for the community of researchers, only until NASA released a large 
evidential dataset in 2004~\citep{Sayyad-Shirabad+Menzies:2005}.

It it interesting to mention that
in early 2000s most authors assumed that ``distributions'' in both 
training and testing datasets were similar~\citep{Turhan2009} (\emph{intra-release}), which
mislead them because the assumption was wrong.

Later,~\citet{5609530} introduced ``ensemble techniques,'' which made the analysis
possible in situations when software modules do not have similar 
distributions in their training and testing datasets (\emph{inter-release}).

Recently,~\citet{Better_cross} demonstrated that defect detection may work across 
different software modules, using their historical data (\emph{cross-project release}).

\subsection{By the Output}

ML-based defect detection methods can be classified by the type of the output they produce:

First, the output could be \emph{binary}, which means that 
software modules (like classes, methods, packages, etc.) are classified
either as faulty or non-faulty ones, as done by~\citet{Gokhale97regressiontree, Menzies04assessingpredictors}.

Second, \emph{defect density} or the \emph{number of defects} 
may be reported per module, demonstrating the degree of their defectiveness,
as suggested by~\citet{Predicting_Fault, JANES20063711, 7510216}.

Third, \emph{defect severity prediction} may be reported per module, demonstrating
the impact of defects found in the module to the end-user experience,
as suggested by~\citet{7510216, SHATNAWI20081868, 1717471}.

Fourth, the output could contain information about a specific location
of the defect inside the source code (the line of code with a defect), 
as suggested by~\citet{vasic2019neural}.

\subsection{By the Input}

\textbf{Software metrics} can be used as an input. 
They can be obtained from the source code,
and can belong to one of the following categories:

\begin{itemize}
\item \emph{Product metrics} are generally used to check 
whether software follows~\citep{InternationalStandardOrganization}, 
as it was mentioned  by~\citet{10.1007/s10462-017-9563-5}, 
and can be classified as suggested by~\citet{10.5555/540137}: 
\begin{enumerate*}[label=\arabic*)]
\item \emph{Traditional} metrics may include size, 
system complexity~\citep{1702388}, and others;
\item \emph{Object-oriented} metrics may include coupling, 
cohesion and others specifically related to object-oriented 
methodology, as suggested by~\citet{979986, Incorporating_transitive};
\item \emph{Dynamic} metrics are gathered from a running  
program and demonstrate the behavior of a software component
during its execution~\citep{MITCHELL20064}.
\end{enumerate*}

\item \emph{Process metrics} such as number of modules changed for a bug-fix,
 work products delivered, and so on may also be used
 as an input~\citep{The_IT_Measurement_Compendium, 10.1109/ISSRE.2010.25}.
\end{itemize}

\textbf{AST} can also be used
as an input, which is a represention of the syntax of a 
software code snippet as a tree-like structure. 
E.g.,~\citet{6676914} computes the set of AST changes between two source
code files using algorithm, published by~\citet{4339230}.

\subsection{By the ML Model}

There are plenty of studies for bug detection problem 
in which ML methods are used. ML methods 
for solving defect detection problem can be classified into the following groups: 

\begin{enumerate}[label=\arabic*)]
\item \emph{Unsupervised learning} is a learning model which uses unlabeled data.
The goal of unsupervised methods is to uncover unknown patterns in the data. 
Unsupervised learning for solving defect detection problem includes such algorithms 
as K-Means clustering cleaning approach~\citep{8952192}, KNN~\citep{8777507}.

\item \emph{Supervised learning} uses labeled data to produces a function, 
which maps unseen data to labels. Supervised methods for defect detection includes
\citep{10.1145/1137983.1138012}, Naive Bayes~\citep{4027145}, Random forest 
\citep{1383136, 7476673}, nearest neighbor 
\citep{10.1145/2786805.2786813}, SVM~\citep{ELISH2008649, 10_1007}, 
neural network~\citep{THWIN2005147THWIN2005147},~\citep{1033229, 10.1145/3360588} and ensemble techniques~\citep{Ensemble_Techniques}.
Some authors tried to combine different machine learning techniques 
(Linear Regression, Neural Network for continuous goal field, Naive Bayes, etc.) 
with statistical techniques, such as PCA~\citep{1544801}.

\item \emph{Semi-supervised learning} uses both labeled and unlabeled data and
is a combination of supervised and unsupervised methods. 
It can help unsupervised methods to obtain better-defined clusters,
as explained by~\citet{semi_supervised_learning, 7965301, Zhang2017}.

\item \emph{Deep learning} is a class of ML algorithms, which extracts
features from lower layers in order to use them in higher ones.
Textual information (for example, English bug reports
or source files comments and names of variables) can be used in deep learning. 
Using this information, deep learning 
can achieve better performance than many state-of-the-art approaches.
For example,~\citet{8616596} proposed a different method based 
on Long Short-Term Memory (LSTM) networks 
and~\citet{XIAO201917} suggested to use 
convolutional Neural Network together 
with word-embedding and feature-detecting techniques. 
\end{enumerate}


\section{Available data}
\label{sec:data}
We considered available data in some open sources 
related to the defect detection problem. Some datasets 
have only code samples and labels, while others have 
calculated some code metrics.

ETH JavaScript Dataset\footnote{\url{https://www.sri.inf.ethz.ch/js150}}
includes 150,000 JavaScript code samples represented as AST trees. 
The dataset has been used in a binary classification task. 
Target positive examples were extracted from the code 
and negative examples were generated using 
Negative Sampling techique. The dataset 
was used by~\citet{DeepBugs} in an attempt to find the following defects: 
swapped arguments, wrong binary operator, and wrong binary operand.

NASA Dataset\footnote{\url{http://promise.site.uottawa.ca/SERepository/datasets-page.html}}
includes 20 small-sized datasets (20 to 2000 samples in each one) of NASA projects. 
It does not contain code itself but only code metrics. The data 
was previously used for defect prediction, cost estimation, successful reuse, and requirements tracing.
 
Defects4j Dataset\footnote{\url{https://github.com/rjust/defects4j/tree/additional-projects-4}}
includes code samples from different projects, 
where each one has 52 defects of different 
type. The datasets was built by scrapping bug reports and commits 
that fixed them. This dataset does not have labels.

BugLocalization Dataset:\footnote{\url{https://github.com/yanxiao6/BugLocalization-dataset/tree/master/dataset}}
has 18,500 Java defect report samples from five GitHub-hosted projects. 
Each sample consists of defect description, a commit to fix the defect 
and a link to GitHub source code file.

CodeSearchNet dataset:\footnote{\url{https://github.com/github/CodeSearchNet}}
is a large dataset of over two million rows of code 
collected from open source libraries. The dataset contains 
Python, Javascript, Ruby, Go, Java, and PHP languages. 
Each sample consists of a function code and a comment related 
to the function. The dataset is built for a semantic code search task.


\section{Selection Criteria}
\label{sec:criteria}
There is a number of crucial expectations we have for the
method and the tool under design:

\textbf{Novelty} is crucial and has a few important aspects:
\begin{enumerate*}
\item 
Novelty of the \emph{method} is one the most important values of novelty. 
E.g.,~\citet{Akiyama1971AnEO} was the first who 
tried to solve defect detection problem using 
linear regression as it was mentioned by~\citet{7476771}.
Sometimes authors can use some new idea in the 
algorithm, e.g,~\citet{XIAO201917} used convolutional neural network
together with word-embedding and feature-detecting techniques.

\item
The \emph{effectiveness} of the model is another important 
value of novelty. If authors published new model which 
is more effective, it means that we can solve a problem 
much better. All evaluation measures can be described later.

\item
Collecting new dataset is one of the most important and hard work in Data Science. 
Main results can vary depending on method used when gathering a dataset. 
It is necessary to do it in a correct way, since there are many issues related 
to data problems like \emph{Outliers}, \emph{Class Imbalance}, \emph{Data shift Problem}, 
\emph{High Dimensionality of Data}, etc. mentioned by~\citet{10.1007/s10462-017-9563-5}.
Also, \emph{granularity} is another problem in the 
defect detection problem. Defect can be found in a module, 
function, line of the code, etc. That is why seems 
it is not an easy task to compare the models with a
different level of granularity.
E.g,~\citet{6464273} analyzed five papers published by 
IEEE and mentioned that there are some data 
quality issues in these papers.
\end{enumerate*}

\textbf{Evaluation} is the essential of ML. It is impossible 
to say whether a model is good or bad, not evaluating it. 
A lot of model performance evaluation measures were used 
for defects detection, but we can 
classify them the following way \citep{10.1007/s10462-017-9563-5, Jiang2008}:
\begin{enumerate*}[label=\arabic*)]
\item \emph{Numeric} performance evaluation measures
are mostly used for defect detection models~\citep{Jiang2008}. 
They can include accuracy, F1-score, G-means, specificity, 
f-measure, and so on. E.g., they are used by~\citet{6349519}. 
\item \emph{Graphical} measures are graphs derived from 
the confusion matrix~\citep{Jiang2008}. 
They can include ROC curve, precision-recall curve, cost curve etc. 
E.g., they are used by~\citet{SHATNAWI20081868}.
\end{enumerate*}

\textbf{Performance} is another important value for defect detection model. 
The selection of the best performance evaluation measure 
is not a trivial task.
For example,~\citet{Jiang2008} compared different alternatives
and demonstrated that no single performance evaluation measure is able to evaluate the
performance of a defect prediction model fair enough. The authors also added 
that it is better not only to measure model classification performance 
(like accuracy, recall, etc.), but to minimize the misclassification cost.
\citet{ARISHOLM20102} also compared different types of measures 
and mentioned that it is hard to draw general conclusions about the 
best performance evaluation measure.

Since the tool will be actively used with Continuous Integration (CI), 
it must have a reasonable performance: it must predict in a few seconds. 
If calculation takes a lot of time, a developer may stop using it.
\citet{humble2010continuous} also mentioned that speed 
is very important for CI. The sooner you release the software, 
the sooner you get a return on your investment. 
Performance is also important because it is necessary 
to know whether bugfixes are useful. That is why
we need to minimize the delay between the releases and 
thus, accelerate the feedback, as explained by~\citet{humble2010continuous}.

State-of-the-art papers use Neural Networks (NN) for defect 
detection problem since it can give better results~\citep{XIAO201917,10.1145/3360588}. 
Training takes a lot of time and demands TPU or GPU resources if we use NN. 
Word embedding and other feature-detecting techniques can only increase training time. 
E.g., \citet{10.1145/3360588, 8616596} noticed that 
the training time matters for defect detection problem and 
demonstrates the training time of their models.



\section{The Method}
\label{sec:method}
%Here will be paragraph linking previous sections with "Code Readability".

%The research question is the following: is it possible, using static code analysis to do 
%a code refactoring recommendations to improve a code readability? 

We are looking for an answer to the following research question: 
Is it possible,
combining static code analysis and ML, 
to detect defects in a Java class and 
to give specific recommendations for its refactoring? 

To answer this question we are doing the following research:
\begin{enumerate*}[label={\alph*)}]
    \item We take a set of Java \emph{classes};
    \item We collect a number of static analysis \emph{metrics} per each class;
    \item We locate \emph{code patterns} inside each class;
    \item We ask volunteer programmers to review 
      each class to estimate their \emph{readability};
    \item We put collected data together into a \emph{dataset}; 
    \item We find \emph{relations} between metrics and patterns in the dataset; 
    \item We take an unseen Java class and \emph{locate code patterns}, which impact the quality more then others.
\end{enumerate*}
    
\textbf{Classes}.
We parse Java projects from GitHub to obtain training Java classes. 
All parsed projects must have more than 100 stars, and more than five collaborators. 
Collected Java classes must have more than 50 lines of code and less than 300.  

\textbf{Metrics}.
The code metrics are calculated using existing open-source
tools like CheckStyle and SourceMeter. Examples of metrics are Lines of Code (LoC), 
Cyclomatic Complexity (CC), and Number of Incoming Invocations (NII).

\textbf{Code Patterns}.
AST patterns are the features that reflect the presence 
of some syntax structures in the code, 
for example number of nested \texttt{FOR} loops 
or the amount of ternary operators. The AST patterns are manually designed 
and always point to a specific place in the code. 

\textbf{Readability}.
To gather readability characteristics we conduct
a survey, where volunteering programmers are asked to estimate
the readability of training Java classes, giving answers on a $[0..9]$ scale, 
where 0 means not readable at all and 9 means perfectly readable.
Each programmer is asked to review some snippets from 
the entire training dataset, which means that each
snippet is reviewed by a few programmers. The readability ``score'' per
snippet is a mean of all answers collected.

\textbf{Dataset}.
We combine and put the data into the single dataset. 
In the \autoref{tab:table1} you can see the example how dataset looks
(here CC, CBO, LCOM, and NMD are acronyms for software metrics, 
while RS is the Readability Score collected from volunteers).

\begin{table}[H]
\begin{center}
\begin{tabular}{rrrrrr}
\hline
CC & CBO & LCOM & NMD & \dots & RS \\
\hline
3 & 6 & 34 & 2 & \dots & 4.3 \\
4 & 5 & 55 & 1 & \dots & 2.4 \\
3 & 5 & 22 & 0 & \dots & 5.2 \\ 
\hline
\end{tabular}
\end{center}
\caption{Example of collected dataset}
\label{tab:table1}
\end{table}

\textbf{Relations}.
Using ML methods we are going to learn how code metrics, 
AST patterns and readability are related. 
We are planning to use ML techniques to investigate the importance of features 
with respect to the readability and find combinations that have stronger impact.

\textbf{Locate Code Patterns}.
Finally, for an unseen Java class we can calculate all features except the readability. 
Using known features and the knowledge about relations between features we are planning
to give refactoring recommendations and provide links to the lines of code 
where most problematic patterns are located.   



\section{Risks}
\label{sec:risks}
There is a number of risks we identified and expect to mitigate.

\textbf{Survey Validity}.
While designing the survey procedure we have to 
consider the following threats to its validity:
\begin{enumerate*}[label=\arabic*)]
    \item Set of interviewers must be representive. 
    Will the results be different if we take a different set of interviewers?
    \item Set of code snippets must be representive. 
    Will the results be different if we take a different set of code snippets? 
    \item How to identify bad interviewers? 
    Some interviewers may give answers that are not correlated with actual readability.
\end{enumerate*}
To mitigate this risk we have to think about the diversity of interviewers, 
for example by age, skills, experience, education, and so on. We have to select code snippets
that vary in length and project domain. The more data we collect the better. 
We have to preprocess survey results to exclude outliers.

\textbf{Feature Limits}.
To predict the readability we are going to use features 
calculated with static code analysis tools. 
These features mainly reflect the structure, the syntax and 
size properties of a particular code snippet.  
However, we are not considering semantic properties of the code. 
In the future, we can design semantic related features and 
add them to our dataset.

% \section{Implementation Details}
% \label{sec:implementation}
% \input{sections/implementation}

\section{Empirical Results}
\label{sec:results}
In our first experiment we collected dataset of Java classes and calculated 
metrics and patterns for them. We calculated the readability using open-source
tool. Then we estimate the impact features to readability and conclude that
\emph{ncss\_method\_avg}, \emph{halstead\_volume}, and \emph{right\_spaces\_var} are
most important features. The more details you can find in \emph{Report 1}.

In second experiment we figured out non-zero correlation between the 
non-structural patterns and the code complexity metrics like \emph{Cyclomatic Complexity}
and \emph{NPath}. We observe that \emph{var\_in\_the\_middle} shows  
highest impact on complexity. We explain it by the fact that the more complex 
programs will  tend to have more variable and some of them declared in the middle
of the method body. The more details you can find in \emph{Report 2}.



\section{Conclusion}
We will write this section later.

\section{Future Work}
We will write this later.

\section{Acknowledgements}
The tool was designed in System Programming Lab. Many thanks
to the contribution of (in alphabetic order of the last name):
Yegor Bugayenko,
Anton Cheshkov,
Vadim Chibiriev,
Lu Jianhua,
Evgeny Maslov,
Yu Jiayuan,
Alexey Zorchenkov.

%\section{Notes}
%Code readability is closely related to defect detection problem. There are plenty of 
studies where the problem of code readability is considered 
\citep{8651396, xxx66666444, 10.1145/1985441.1985454}. At the moment, it is 
still the issue since the problem is too subjective. 

It is interesting that \citet{10.1109/ICPC.2019.00014} tried to evaluate 
the models for code readbility problem. The authors extracted 548 commits 
from 63 engineered Java software projects. The authors identified a commit as a readability commit
when the message of the commit directly indicated it. E.g., 
the authors searched such words of the commit as ``readable'',
``readability'', ``easier to read'', ``understand'', etc.
Several state of the art readability models were reviewed by \citet{10.1109/ICPC.2019.00014} 
and it was discovered, that the models failed to capture readability improvements. 
The authors also suggested several metrics which were not considered in the reviewed state-of-art papers.
The authors believe that those metrics for code readability can be used to detect readabilty 
changes more efficiently.

It is important to notice that code readabilty also can have a correlation 
with complexity of the code. \citet{10.1109/TSE.2009.70} compared code readability 
and software complexity. They used tools, published by \citet{xxx66666444, Readabil74:online}, to
compute metrics for code readability like \emph{The Automated Readability Index}, 
\emph{The Simple Measure of Gobbledygook}, \emph{Flesch-Kincaid Readability Index}
\emph{The Gunning's Fog Index}, \emph{Coleman-Liau Index and Buse Readability Score}. 
Also, they computed complexity software metrics, such as \emph{Halstead Complexity Volume},
\emph{McCabe's Cyclomatic Complexity}. The authors found out that there was a 
negative correlation between the readability and complexity metrics.
It means that low readability increases program complexity and vice versa. 
The authors also mentioned that the languages constructions as comments, spacing, 
while loop, meaningful names and do-while loop affects the code readiblity the most.
The authors also published a dataset with Java code which can help to detect defects.

\citet{10.1007/978-3-319-95171-3_32} proposed a model which quantitatively measures the readability 
of source code. The first idea is the using metrics (\emph{LOC, ProgramVolume, Entropy})
as the key indicators which affect the source code readability. The second
idea is that the authors introduced the equation for quantitative measure 
the source code readability in real time. This measure can demonstrate 
how small changes in the code affect the readability of the code.
The authors also tried to optimize the model and they reached about 
74.59\% of explanatory power. 


\bibliography{references}




\end{document}
